\chapter{Cơ sở lý thuyết}
\label{chap:theory}

Chương này tổng hợp các khái niệm, mô hình và công nghệ được sử dụng. Công thức hoặc ngưỡng do đề tài xây dựng được chuyển sang thiết kế chi tiết ở Chương~\ref{chap:results}.

\section{Dữ liệu tài chính và tiền xử lý}
\label{sec:fin-data}

\subsection{Sổ cái quan hệ và toàn vẹn dữ liệu}

Một giao dịch tài chính tối thiểu gồm số tiền, loại, ngày, mô tả, danh mục và ví. Thu nhập làm tăng số dư; chi phí làm giảm số dư; chuyển ví tạo hai biến thiên trái dấu nên không được tính thành thu hoặc chi. Các thay đổi liên quan nhiều bản ghi phải nằm trong cùng transaction để hoặc hoàn thành toàn bộ, hoặc được rollback toàn bộ.

Cơ sở dữ liệu quan hệ phù hợp với bài toán vì hỗ trợ khóa ngoại, kiểu số thập phân, ràng buộc miền và truy vấn tổng hợp. Dữ liệu nghiệp vụ bền vững cần tách khỏi trạng thái ngắn hạn như câu hỏi đang chờ, bản xem trước, cache hoặc queue. Phân tích theo ngày/tuần/tháng phải giữ trục lịch và điền 0 cho kỳ không phát sinh; bỏ các kỳ 0 có thể làm sai slope, burn rate và tương quan.

\subsection{Chuẩn hóa và độ tương đồng văn bản}
\label{subsec:text-similarity}

Đầu vào được chuẩn hóa chữ thường, dấu câu, khoảng trắng, đơn vị tiền và ngày tương đối trước khi ánh xạ sang schema. Với hai chuỗi $a,b$, độ tương đồng dựa trên khoảng cách Levenshtein \cite{levenshtein1966} được dùng để đo số phép sửa ký tự:

\begin{equation}
  s_{edit}=1-\frac{D_{lev}(a,b)}{\max(|a|,|b|)}.
\end{equation}

Với hai tập token $T_a,T_b$, hệ số Dice \cite{dice1945} là

\begin{equation}
  s_{dice}=\frac{2|T_a\cap T_b|}{|T_a|+|T_b|}.
\end{equation}

Hai thước đo bổ sung cho nhau: Levenshtein phù hợp lỗi gõ gần nhau, còn Dice phù hợp cụm từ có thứ tự hoặc thành phần khác nhau. PERFIN kết hợp exact match, alias, hai điểm tương đồng và một khoảng cách an toàn giữa hai ứng viên tốt nhất. Ngưỡng cụ thể là lựa chọn thiết kế, được trình bày ở mục~\ref{subsec:classification-design}.

\section{Các kỹ thuật phân tích giải thích được}
\label{sec:algorithms}

\subsection{Xu hướng, bất thường và tương quan}

Hồi quy tuyến tính bình phương tối thiểu (OLS) \cite{montgomery2012} mô tả xu hướng của chuỗi theo thời gian bằng $\hat y_i=a+bx_i$. Hệ số dốc

\begin{equation}
  b=\frac{\sum_i(x_i-\bar{x})(y_i-\bar{y})}
          {\sum_i(x_i-\bar{x})^2}
\end{equation}

cho biết chiều biến đổi; $R^2$ cho biết mức độ đường thẳng giải thích dữ liệu. OLS dễ diễn giải và kiểm thử, nhưng không mô hình hóa mùa vụ hoặc sự kiện tương lai.

Z-score so một quan sát với trung bình theo đơn vị độ lệch chuẩn, còn khoảng tứ phân vị (IQR) bền vững hơn trước giá trị cực đoan \cite{tukey1977}. PERFIN dùng hai kỹ thuật như tín hiệu yêu cầu người dùng rà soát, không dùng để kết luận gian lận.

Hệ số Pearson \cite{pearson1895} đo mức đồng biến tuyến tính của hai chuỗi:

\begin{equation}
  r=\frac{\sum_i(x_i-\bar{x})(y_i-\bar{y})}
  {\sqrt{\sum_i(x_i-\bar{x})^2\sum_i(y_i-\bar{y})^2}}.
\end{equation}

Tương quan chỉ được diễn giải trong cửa sổ quan sát, không hàm ý quan hệ nhân quả.

\subsection{Runway, khoản định kỳ, ngân sách và mục tiêu}

Runway là ngoại suy số ngày số dư hiện tại có thể duy trì nếu tốc độ chi gần đây không đổi. Khai phá khoản định kỳ nhóm giao dịch theo mô tả, mức tiền và khoảng cách ngày. Dự báo ngân sách dùng tốc độ chi theo số ngày đã trôi qua; lập kế hoạch mục tiêu dùng phần tiền còn thiếu, hạn chót, mức góp và mô phỏng trả nợ. Các kỹ thuật này ưu tiên khả năng giải thích và xử lý điều kiện biên hơn độ phức tạp mô hình. Công thức, cửa sổ và quy tắc do PERFIN lựa chọn được trình bày ở mục~\ref{subsec:analytics-design}.

\begin{table}[H]
  \centering
  \caption{Kỹ thuật phân tích và lý do lựa chọn}
  \label{tab:algorithm-selection}
  \begin{tabularx}{\textwidth}{@{}L{3.0cm}Y Y@{}}
    \toprule
    \textbf{Kỹ thuật} & \textbf{Vai trò} & \textbf{Lý do phù hợp} \\
    \midrule
    OLS & Xu hướng và dự báo ngắn hạn. & Ít tham số, giải thích được slope và mức phù hợp. \\
    Z-score + IQR & Phát hiện khoản/ngày chi bất thường. & Kết hợp độ nhạy theo phân phối và độ bền vững theo quantile. \\
    Runway & Cảnh báo nguy cơ cạn dòng tiền. & Phép tính trực tiếp, dễ nêu giả định. \\
    Gom cụm theo quy tắc & Nhận diện khoản chi lặp. & Phù hợp tập dữ liệu nhỏ, không cần huấn luyện. \\
    Pearson & Tìm cặp danh mục đồng biến. & Có hệ số chuẩn hóa và cách diễn giải rõ. \\
    Budget/Goal Planner & Dự báo hạn mức và tiến độ. & Kết quả có thể tính tay và mô phỏng what-if. \\
    \bottomrule
  \end{tabularx}
\end{table}

\section{OCR, STT và mô hình ngôn ngữ}

\subsection{OCR và Speech-to-Text}
\label{subsec:ocr-stt}

OCR thường gồm tiền xử lý ảnh, phát hiện vùng chữ, nhận dạng và hậu xử lý; Tesseract là một kiến trúc OCR điển hình được mô tả trong \cite{smith2007}. STT chuyển tín hiệu âm thanh thành transcript; Whisper minh họa hướng nhận dạng giọng nói học từ dữ liệu quy mô lớn \cite{radford2022}. Trong PERFIN, cả OCR và STT chỉ tạo raw text kèm tên provider. Raw text tiếp tục qua pipeline chuẩn hóa, xem trước và xác nhận như dữ liệu nhập tay.

Chất lượng OCR/STT phải được đo bằng ground truth. CER/WER đo sai số văn bản, còn transaction-field accuracy đo tác động lên số tiền, ngày, danh mục hoặc dòng hàng. Một smoke test với vài tệp chỉ chứng minh provider có thể chạy, không phải phép đo accuracy.

\subsection{Parser cục bộ và LLM}

Transformer dùng attention để mô hình hóa quan hệ giữa token và là nền tảng của LLM hiện đại \cite{vaswani2017}. Mô hình đa phương thức như Gemini mở rộng khả năng hiểu văn bản và media \cite{gemini2023}. Tuy nhiên, structured output chỉ bảo đảm hình dạng dữ liệu; nó không bảo đảm số tiền đúng, danh mục thuộc người dùng hay một thao tác được phép.

\begin{table}[H]
  \centering
  \caption{So sánh parser cục bộ và LLM}
  \label{tab:parser-llm-comparison}
  \begin{tabularx}{\textwidth}{@{}L{3.0cm}Y Y@{}}
    \toprule
    \textbf{Tiêu chí} & \textbf{Parser cục bộ} & \textbf{LLM qua function calling} \\
    \midrule
    Độ bao phủ ngôn ngữ & Tốt với mẫu đã biết, yếu với câu tự do. & Hiểu biến thể và ngữ cảnh tốt hơn. \\
    Tính tái lập & Cao, chạy offline và không tốn API. & Phụ thuộc model/provider và có độ trễ. \\
    Kiểm soát & Quy tắc và lỗi dễ truy vết. & Phải khóa schema, validation và fallback. \\
    Vai trò chọn trong PERFIN & Baseline, fallback và quality gate. & Trích xuất/định tuyến chính khi được cấu hình. \\
    Quyền nghiệp vụ & Không ghi trực tiếp; trả typed draft. & Không ghi trực tiếp; chỉ đề xuất tool và đối số. \\
    \bottomrule
  \end{tabularx}
\end{table}

PERFIN dùng function calling theo schema của nhà cung cấp \cite{geminifunctioncalling}. LLM đề xuất tool và đối số; backend mới validation, dựng preview và chờ xác nhận. Khi diễn giải insight, LLM chỉ nhận facts đã tính. Cách tách này giữ lợi ích giao tiếp nhưng không biến mô hình thành nguồn sự thật.

\section{Công nghệ và lý do lựa chọn}
\label{sec:tech-rationale}

\begin{longtable}{@{}L{3.0cm}L{5.6cm}L{5.3cm}@{}}
  \caption{Công nghệ, vai trò và phương án thay thế}\label{tab:technology-rationale}\\
  \toprule
  \textbf{Công nghệ} & \textbf{Lý do lựa chọn} & \textbf{Giới hạn/phương án khác} \\
  \midrule
  \endfirsthead
  \toprule
  \textbf{Công nghệ} & \textbf{Lý do lựa chọn} & \textbf{Giới hạn/phương án khác} \\
  \midrule
  \endhead
  React Native + Expo & Một mã nguồn cho form, chat, ảnh, audio và preview trên mobile. & Native riêng chỉ cần khi có yêu cầu thiết bị đặc thù. \\
  Node.js + Express & Phù hợp API hướng I/O, dùng chung JavaScript với hàm nghiệp vụ và test. & Modular monolith đủ cho quy mô đề tài; microservice làm tăng chi phí vận hành. \\
  PostgreSQL \cite{postgresql2026} & Khóa ngoại, \code{DECIMAL}, transaction và truy vấn tổng hợp phù hợp sổ cái. & File/NoSQL khó cung cấp cùng mức ràng buộc và rollback. \\
  Redis \cite{redisexpire} & TTL cho pending/clarification, cache và dữ liệu queue. & Không làm nguồn dữ liệu chuẩn; memory fallback chỉ dùng phát triển. \\
  BullMQ \cite{bullmq2026} & Lịch chạy, retry và job ID cho tác vụ idempotent. & Cron đơn giản không có cùng cơ chế queue/retry. \\
  Gemini + local parser & Kết hợp độ bao phủ ngôn ngữ với fallback kiểm thử được. & Provider nằm sau adapter để không khóa logic nghiệp vụ. \\
  OCR/STT cục bộ hoặc cloud & PaddleOCR/PhoWhisper ưu tiên riêng tư; Google Cloud tăng khả năng thay thế provider. & Mọi raw output vẫn phải qua cùng validation/preview. \\
  \bottomrule
\end{longtable}

\section{Ứng dụng liên quan và hướng lựa chọn}

Money Lover \cite{moneylover} và MISA MoneyKeeper \cite{misamoneykeeper} đại diện cho nhóm ứng dụng quản lý tài chính dựa trên biểu mẫu và dashboard. Nhóm này mạnh về dữ liệu cấu trúc và báo cáo nhưng vẫn cần nhiều thao tác nhập. Chatbot giảm thao tác nhưng có rủi ro trả số không truy vết được; hệ thống phân tích chuyên sâu lại khó dùng với người dùng phổ thông.

PERFIN không đặt mục tiêu cạnh tranh về độ đầy đủ thương mại. Khoảng trống được chọn là kết hợp bốn yếu tố: sổ cái quan hệ, giải thuật xác định, con người trong vòng xác nhận và lớp ngôn ngữ có ranh giới. Lựa chọn này phù hợp trọng tâm dữ liệu--giải thuật của niên luận và tạo điều kiện đánh giá từng thành phần độc lập.

\begin{table}[H]
  \centering
  \caption{Tổng hợp khoảng trống và giải pháp của PERFIN}
  \label{tab:research-gaps}
  \begin{tabularx}{\textwidth}{@{}Y Y Y@{}}
    \toprule
    \textbf{Khoảng trống} & \textbf{Rủi ro} & \textbf{Giải pháp được chọn} \\
    \midrule
    Nhập tự nhiên dễ ghi sai & Dữ liệu bẩn lan sang báo cáo & Typed draft, validation, preview và confirm. \\
    Media chỉ cung cấp raw text & Sai số tiền/ngày/danh mục & Ground truth, đo field accuracy và xác nhận. \\
    Chatbot trả số không rõ nguồn & Hallucination, khó kiểm toán & SQL/analytics tạo facts; narrator chỉ diễn giải. \\
    Phân loại không thích nghi & Lặp lại cùng lỗi & Correction retrieval có ngưỡng và kiểm soát. \\
    Retry tạo hiệu ứng trùng & Sai dữ liệu hoặc thông báo & Job ID, fingerprint và thao tác idempotent. \\
    \bottomrule
  \end{tabularx}
\end{table}
